\section{Introduction}

Four-bit floating-point (FP4) Tensor Cores accelerate general matrix
multiplication (GEMM), but each training layer must also produce scales, pack
values, build forward and backward layouts, and retain autograd state. When
separate kernels move bfloat16 (BF16) intermediates through high-bandwidth
memory (HBM), this surrounding work can erase the GEMM speedup
\cite{nvidia2025nvfp4}.

We address this gap with \emph{format-aware fusion}: each producer completes
the scale shared by its output values while emitting the packed layouts and
saved state needed by the following matrix multiplications. We use
\emph{route} to mean a complete numerical recipe and execution path, not only
an FP4 datatype. \mxfp{} completes a local 32-value scale, global
\nvfp{} waits for a tensor-wide maximum, and cooperative-thread-array-local
(CTA-local) \nvfp{} completes two outer scales inside the tile that owns the
matrix multiplication. The profitable fusion boundary therefore differs by
format.

We evaluate the resulting paths in Llama-3-family 8B pretraining at sequence length
8,192 and global batch 512. Same-accelerator probes measure every route with a
BF16 output projection and compiled cross entropy (CE), while long
trajectories cover up to 160B tokens. The fastest custom route reaches 37.9K
tokens/s/GPU. The \mxfp{} route with row-gradient stochastic rounding (SR)
and fixed-sign 32-value Hadamard (H32) weight-gradient (Wgrad)
preconditioning reaches 37.2K
tokens/s/GPU and ends 2.11\% above the raw BF16 endpoint. The TE-native and
four-final-BF16-block trajectories end 1.34\% and 0.87\% above raw BF16,
respectively, at measured throughputs of 27.6K and 27.1K tokens/s/GPU.

Table~\ref{tab:prior-work-boundary} situates the implementation relative to
established FP4 practice. We contribute:

\begin{itemize}
  \item \emph{Numerical parameterization.} We specify the FP4 contract per
        operand: upward-rounded \mxfp{} scales for activations and gradients;
        an adaptive finer scale for $32\times32$ weight tiles; separate FP32
        row- and column-tile scale factors for \localcta{}; stochastic
        rounding only for the data-gradient multiplication; and paired
        fixed-sign 16- or 32-value Hadamard preconditioning only for the
        weight-gradient multiplication.

  \item \emph{Format-aware fusion and native implementation.} We implement
        custom \mxfp{}, custom global \nvfp{}, and \localcta{} producers that
        finish the scale calculation over the required group of values and emit the layouts and saved
        state required by forward and backward. The MXFP4 and CTA-local weight
        producers quantize once and emit both packed weight layouts in the same
        fused operation. Custom global NVFP4 likewise integrates its shared
        $16\times16$ weight encoding and both packed layouts directly into the
        fused path. Backward stages fuse optional Hadamard preconditioning
        with column quantization. The same framework supports depth- and
        operand-wise hybrids.

  \item \emph{End-to-end empirical evaluation.} We implement and evaluate the
        custom pure and hybrid routes alongside BF16 and TE comparators. The
        hybrid designs include depth-wise format assignment and an operand-wise
        assignment motivated by a data-gradient error analysis using identical
        saved tensors. The evidence
        comprises matched throughput probes, trajectories through 160B
        training tokens, checkpoint-matched validation loss on a fixed
        external data stream, and
        downstream evaluation of the complete recipes. We release the
        training and evaluation code, pinned kernel sources, complete recipe
        ledger, publication data, and figure builders at
        \url{https://github.com/MrHuff/mfu-fp4-training};
        Appendix~\ref{app:code-release} describes its reproducibility contract.
\end{itemize}
